\section{The outgoing collar, the protected disk and terminal Ricci charge}
\label{sec:ed03-collar-charge}

\begin{theorem}[Comparing the retained and calibrated collars]
\label{thm:ed03-actual-collar-transition}
\uses{thm:ed03-original-rounded-trace}
\lean{SmoothPoincare.IsolatedOneHandle.RetainedRoundedTrace.exists_outgoingCollarComparison, SmoothPoincare.IsolatedOneHandle.RetainedRoundedTrace.exists_outgoingCollarComparison_toNextCut, SmoothPoincare.IsolatedOneHandle.outgoingCollarTransition_fixedOn_iff, SmoothPoincare.IsolatedOneHandle.outgoingCollarTransition_normalDerivative_pos}
\leanok
Let $A$ be the calibrated source chart, $F$ the two-sided extension of
the retained inward map, and $P$ the target product parametrized by the
actual boundary marking. There is a smooth coordinate change
$J=P^{-1}F$ fixing the entire zero section. Its domain and inverse domain
contain the same full strip of positive width. Both $J$ and $J^{-1}$
preserve the inward and outward sides. Their tangential derivatives at the zero section are the identity and their normal derivatives there are strictly positive. At calibrated parameter $q=(x,t)$
the source chart becomes $A(J^{-1}q)$. Its rounded height is
$c_{\rm round}+(J^{-1}q)_2$. The actual target height is $c+t$
and the corresponding next cut has height $h_{\rm next}+t$.
On the inward side the retained body map and inverse give these actual
points. A prescribed parameter set is fixed by the transition precisely
where the two parametrized collars agree.
\end{theorem}
\begin{proof}
\leanok
Compose the partial diffeomorphisms on their open overlap and use
compactness of the full boundary level to put a two-sided strip in both
coordinate domains. The inward map takes strict interior to strict
interior. The calibration and injectivity separate the opposite side.
Apply the two inverse equations to each parameter of a framed tube.
\end{proof}

\begin{theorem}[Covering an independently prescribed ball]
\label{thm:ed03-target-ball-cover}
\uses{thm:ed03-marked-glue,thm:ed03-birth-survivors}
\lean{SmoothPoincare.exists_markedPatchDiffeomorph_of_unitBallCover, SmoothPoincare.IsolatedOneHandle.connectedTrace_of_twoBirthCharts, SmoothPoincare.IsolatedOneHandle.RetainedRoundedTrace.exists_markedBall_of_nativeRepCover, SmoothPoincare.markedEdgeModelFlatten_images, SmoothPoincare.markedEdgeModelCollar_coordinates, SmoothPoincare.markedEdgeRoundBallCollar_coordinates}
\leanok
On a binary merger the two connected incoming body images are joined by
the handle at its two feet. Suppose four local diffeomorphisms to the
independently specified closed four-ball agree forward on source overlaps,
and cover the images of every lower and handle representative in the
rounded body. They cover the entire target, glue to a diffeomorphism and
give its inverse on every patch target. One explicit Euclidean model is
$q(u,z)=(u^2-1)^2+|z|^2$. For $a>1$ the change
\[
 (u,z)\longmapsto
 \left(\frac{u}{\sqrt{a+1}},
      \frac{z}{\sqrt{(a+1)(a-1+u^2)}}\right)
\]
carries $q\le a^2$ onto the unit ball and its full level onto the unit sphere.
Writing this map as $d_a$, a collar of the round target is
$d_a d_{\sqrt{a^2+t}}^{-1}(v)$, for $v$ on the unit sphere and
$|t|<w<a^2-1$. It fixes every $v$ at zero. It lies in the closed ball
precisely for $t\le0$, with a smooth inverse recovering both parameters.
\end{theorem}
\begin{proof}
\leanok
The native lower and handle representatives exhaust the adjunction.
The glued map is a local diffeomorphism from a compact connected body.
The resulting covering of the simply connected ball has one sheet.
Its inverse agrees with each patch inverse by injectivity. The explicit
model follows by factoring the height defect into a positive function
times the difference of the squared norm from one.
\end{proof}

\begin{proposition}[The closed split and slid controls]
\label{prop:ed03-closed-slide-control}
\uses{thm:ed03-marked-glue}
\lean{SmoothPoincare.ClosedMarkedControl.splitExterior_torus, SmoothPoincare.ClosedMarkedControl.splitTorusCollar_transport, SmoothPoincare.ClosedMarkedControl.exists_splitClosingHomeomorph, SmoothPoincare.ClosedMarkedControl.slid_dualTube_and_inverse, SmoothPoincare.ClosedMarkedControl.slidClosedAssembly_representatives}
\leanok
In the split control the spectator in $S^3\subset\mathbb C^2$ has tube
$(z,v)\mapsto (\sqrt{1-\epsilon^2|v|^2}\,z,\epsilon v)$,
$0<\epsilon<1$. The explicit complementary solid torus agrees with
this tube on both boundary parameters. Its smooth signed collar recovers $t=|y_2|^2-\epsilon^2$ and both circle parameters throughout $-\epsilon^2<t<1-\epsilon^2$. The two marked disk pieces glue to $S^1\times S^2$ by the lower and upper hemisphere maps, with an inverse on each piece. Given the slid trace comparison
$F$, its boundary restriction and its agreement on the outgoing
spectator chart, the upper closing map is $\gamma'=\gamma\circ\partial F$.
The dual framed spectator and its inverse agree with those of the split
control. The full closed adjunction comparison is a homeomorphism,
acting by $F$ below and the identity above, with the inverse formulas
on every representative.
\end{proposition}
\begin{proof}
\leanok
Evaluate the complementary torus chart at its two boundary circle
parameters. For the upper dual chart, swap the two factors and apply
the closing map. Its inverse applies the inverse closing map before
unswapping. The two representative equations descend through the
closed attachment because the closing square commutes.
\end{proof}

\begin{theorem}[A finite nonconformal disk variation]
\label{thm:ed03-signed-first-variation}
\uses{thm:ed03-positive-conformal-core,prop:ed03-curved-cap-price}
\lean{SmoothPoincare.exists_clampedAffineMetricFamily, SmoothPoincare.exists_selectedNonconformalDiskSaving_of_weightedSquare, SmoothPoincare.protectedDisk_positiveCap_allowance_lower}
\leanok
Fix one finite positive source below $2T$, its disk $B$, protected core
$Q$, and native trim $W$. A localized variation supported in the interior
of $B$ away from $Q\cup W$ preserves the prescribed metric germs.
Suppose its signed energy derivative at zero is $-D/2$,
where $D=\int\chi b^2\,dV_g>0$, its tangent at a point of the indicated
affine path is nonscalar, and a continuous sign detector is positive there
at parameter zero. With the stated frame agreement off the support, some
$t>0$ gives a positive nonconformal endpoint. Its saving in each
chirality exceeds $tD/4$ and is at most $e-I$. The whole short joined path
stays strictly below $2T$. If a same-type positive cap costs at most
$\lambda_h^s(V)+\delta_s$, $V=(\mathrm{int}\,B)^c$, retention of disjoint
$Q$ and $B\cap W$ on a subcritical joined endpoint forces
\[
 \delta_s>I+\lambda_g^s(Q)+\lambda_g^s(B\cap W)-2T .
\]
\end{theorem}
\begin{proof}
\leanok
Clamp the real parameter to $[-1/2,1/2]$ for a symmetric direction
with $|D_x(v,v)|\le g_x(v,v)$. Every metric is spatially smooth and
positive definite, and the path is affine near zero. The derivative
and continuity at zero select the finite interval. The signature identity
and the unchanged complete exterior transfer the saving to each signed
disk charge. For the cap bound, add the retained terms inside $B$ and
use $I\le E_s(\mathrm{cap})$ in the cap's own positive class.
\end{proof}

\begin{proposition}[Paying a cap from the net saving]
\label{prop:ed03-net-saving}
\uses{thm:ed03-signed-first-variation}
\lean{SmoothPoincare.diskCharge_add_joinedCost_le_excess_of_cap_paid_from_saving, SmoothPoincare.eventually_no_paid_cap_of_protectedDisk_charge, SmoothPoincare.eventually_no_paid_cap_of_original_charged_markedDisk, SmoothPoincare.original_markedDisk_charge_eq_and_le_excess_of_arbitrarily_small_cap_drift}
\leanok
On an actual finite source let $e$ be the old energy, $h$ the whole
replacement, and $j_s\ge0$ a measured additional joined cost. If the cap
above has allowance $\delta_s$ and $\delta_s+j_s\le e-E_s(h)$,
then $\lambda_h^s(B)+j_s\le e-I$. Along a sequence $e_i\to I$, a
uniformly charged retained core in even one sign, allowed to depend on $i$,
eventually rules out this payment in both signs simultaneously. This
applies when the original marked disk carries a fixed positive charge in
one sign on a selected subsequence and lies inside the protected core.
At one fixed source, suppose for every $\eta>0$ there is a single
positive same-type cap whose energy in each sign differs by less than
$\eta$ from the corresponding original exterior charge. Then the two
original disk charges agree, and each is at most $e-I$.
\end{proposition}
\begin{proof}
\leanok
Use the full identity $E_s(h)=\lambda_h^s(V)+\lambda_h^s(B)$ and the
cap bound. The replacement exterior cancels even when its charge has
changed. Monotonicity on the protected subset and convergence of $e_i-I$
give the sequence conclusion. At one fixed source apply the same-type cap
floor and signature identity at each tolerance before letting it decrease to zero.
\end{proof}

\begin{theorem}[Trace-free Ricci charge in a pointed terminal ball]
\label{thm:ed03-pointed-terminal-charge}
\uses{thm:ed03-terminal-ball-euler}
\lean{SmoothPoincare.tendsto_terminalBallCutoff_weylScalarCharge_of_initial_region, SmoothPoincare.eventually_ricciFlow_terminalBall_intrinsicRicciCharge, SmoothPoincare.intrinsicTraceFreeRicciSq_eq_of_partialDiffeomorph_inner, SmoothPoincare.terminalTraceFreeRicciSq_of_twoJets, SmoothPoincare.pointedTerminalBuffer_tendsto_traceFreeRicci, SmoothPoincare.pointedTerminalBuffer_traceFreeRicci_charge, SmoothPoincare.ricciFlow_terminalBall_intrinsicRicciCharge_pointedLimit}
\leanok
Consider a sequence of smooth modified Ricci solutions with the same fixed
terminal time $\tau>0$, curvature norm bound $A$, integrated derivative
budget $L^2$, and lower Ricci bound $-(b/s)g_s$, $b<1/2$. Suppose each
initial attaining set has full curvature mass at least $\epsilon^2$ and
lies in the radius-$R$ terminal plateau. At time zero the combined
Weyl and scalar-square charge tends to zero on measurable regions
containing the entire radius-$(R+1)$ terminal ball. Fix
\[
 0<c<\frac{\epsilon^2}{2}-288 A L\sqrt{\frac{\tau}{1-2b}} .
\]
Then the terminal ball eventually carries trace-free Ricci mass at least
$c$. Suppose these same terminal metrics and centers have complete pointed
smooth convergence on a cofinal subsequence in the canonical pullback
domains. For every $r>R+1>0$ the limit satisfies
\[
  \int_{B(p,r)} |\mathrm{Ric}^{\circ}|^2\,dV\ \ge c .
\]
\end{theorem}
\begin{proof}
\leanok
Bound the moment of the fixed cutoff by the nonnegative initial charge on
its entire support. Apply the terminal Euler estimate separately on each
solution. On a compact buffer the pulled-back metric two-jets converge.
Coercivity and the jet bounds give pointwise convergence of
$|\mathrm{Ric}|^2-R_{\rm scal}^2/4$ and a uniform spatial bound.
Use the Riemannian volume change of variables and dominated convergence
on that compact set. Choose $R+1<a<r$. The lower metric bound on
$\overline B(p,a)$ prevents a path of source length less than $R+1$ from
exiting its image for all sufficiently large indices. Thus the comparison
of integrals includes the entire source charge ball.
\end{proof}

\begin{proposition}[Volume escape and a positive scalar point]
\label{prop:ed03-volume-and-scalar}
\uses{thm:ed03-pointed-terminal-charge}
\lean{SmoothPoincare.ricciFlow_terminal_sqrtVolume_change_le, SmoothPoincare.ricciFlow_terminal_volume_tendsto_atTop, SmoothPoincare.pointedTerminal_noncompact_of_volume_escape, SmoothPoincare.ricciFlow_terminalPointed_noncompact_of_initialVolume, SmoothPoincare.ricciFlow_exists_scalar_pos_in_charged_ball}
\leanok
For the same connected compact source slices, write $V_i(s)$ for their
total volumes. The spatial curvature budget gives
\[
 |\sqrt{V_i(\tau)}-\sqrt{V_i(0)}|
 \le (\sqrt 6/2)A\tau .
\]
If their initial volumes diverge, their smooth pointed terminal limit
is noncompact. Further suppose the charged limit slice belongs to a
regular limiting Ricci solution. If scalar curvature is nonnegative
on the fixed charged open ball for all times in a neighborhood of the
slice, at a point of that ball both scalar curvature and
$|\mathrm{Ric}^{\circ}|^2$ are strictly positive.
\end{proposition}
\begin{proof}
\leanok
Differentiate $\sqrt{V}$ in the regular interior, use
$|R_{\rm scal}|^2\le6|\mathrm{Rm}|^2$ and Cauchy--Schwarz, then use
continuity at both endpoints. A compact pointed limit would globalize
the comparison and have finite total volume equal to the limit of the
source volumes. Positive integrated charge gives a point of positive
trace-free Ricci square. If its scalar curvature vanished, the time
derivative there would be zero and the spatial Laplacian nonnegative.
The scalar evolution equation has a strictly positive Ricci term there.
\end{proof}

The target-ball theorem still requires native forward charts with the
prescribed overlap and tube equations. The finite disk variation still
requires its energy derivative on the selected physical metric. The
charged terminal limit uses the supplied pointed comparison at positive
time. A comparison at the separated quantiles in the original physical
metric has its own distance, curvature and marking hypotheses.
